\documentclass[12pt,a4paper]{article}

\usepackage[spanish]{babel}
\usepackage[utf8]{inputenc}
\usepackage[T1]{fontenc}
\usepackage{graphicx}
\usepackage{float}
\usepackage{booktabs}
\usepackage{hyperref}
\usepackage{listings}
\usepackage{xcolor}
\usepackage{geometry}
\usepackage{fancyhdr}
\usepackage{titlesec}
\usepackage{enumitem}
\usepackage{tikz}
\usepackage{tcolorbox}

\geometry{margin=2.5cm}

\tcbuselibrary{skins,breakable}

\newtcolorbox{infobox}{
    colback=blue!5!white,
    colframe=blue!75!black,
    title={\textbf{Informaci\'on}},
    fonttitle=\bfseries,
    breakable
}

\newtcolorbox{warnbox}{
    colback=yellow!10!white,
    colframe=yellow!75!black,
    title={\textbf{Precauci\'on}},
    fonttitle=\bfseries,
    breakable
}

\newtcolorbox{tipbox}{
    colback=green!5!white,
    colframe=green!50!black,
    title={\textbf{Consejo}},
    fonttitle=\bfseries,
    breakable
}

\hypersetup{
    colorlinks=true,
    linkcolor=blue,
    urlcolor=blue,
    citecolor=blue
}

\lstset{
    basicstyle=\ttfamily\small,
    breaklines=true,
    frame=single,
    backgroundcolor=\color{gray!10},
    keywordstyle=\color{blue},
    commentstyle=\color{green!60!black},
    stringstyle=\color{red!70!black},
    showstringspaces=false,
    numbers=left,
    numberstyle=\tiny\color{gray},
    tabsize=2,
    language=bash
}

\pagestyle{fancy}
\fancyhf{}
\fancyhead[L]{\leftmark}
\fancyhead[R]{Tutorial --- Monitoreo Solar}
\fancyfoot[C]{\thepage}

\titleformat{\section}{\Large\bfseries}{\thesection}{1em}{}
\titleformat{\subsection}{\large\bfseries}{\thesubsection}{1em}{}

\begin{document}

\begin{titlepage}
    \centering
    \vspace*{2cm}
    {\Huge\bfseries Universidad de Concepci\'on}\\[1cm]
    {\Large Facultad de Ingenier\'ia\\Departamento de Ingenier\'ia El\'ectrica}\\[2cm]

    {\Huge\bfseries Tutorial:\\Sistema de Monitoreo Remoto para Inversor Solar}\\[1cm]
    {\Large Gu\'ia de Uso para Estudiantes y Acad\'emicos}\\[2cm]

    {\large Daniel Ruminot Moscoso}\\[2cm]

    {\large \today}\\[1cm]
    {\large Concepci\'on, Chile}
\end{titlepage}

\tableofcontents
\newpage

% ============================================================
\section{Introducci\'on}
% ============================================================

Este tutorial explica c\'omo usar el sistema de monitoreo remoto del inversor solar Riello H.P.6065REL-D, instalado en el laboratorio de la Universidad de Concepci\'on.

El sistema permite:
\begin{itemize}
    \item Ver datos del inversor en tiempo real desde cualquier navegador
    \item Consultar datos hist\'oricos de energ\'ia, potencia y temperatura
    \item Exportar datos para informes y an\'alisis
    \item Recibir alertas ante condiciones an\'omalos
\end{itemize}

\begin{infobox}
No necesitas instalar ning\'un software. Solo necesitas un navegador web y tu correo @udec.cl.
\end{infobox}

% ============================================================
\section{Acceso al Sistema}
% ============================================================

\subsection{Requisitos}

\begin{itemize}
    \item Navegador web moderno (Chrome, Firefox, Edge, Safari)
    \item Conexi\'on a internet
    \item Correo electr\'onico @udec.cl (o dominios @ing.udec.cl, @dci.udec.cl)
\end{itemize}

\subsection{Primer Acceso}

\begin{enumerate}
    \item Abrir el navegador y dirigirse a:

    \begin{center}
    \url{https://lautaro.elprobedor.com}
    \end{center}

    \item Se mostrar\'a la pantalla de Cloudflare Access pidiendo tu email:

    \begin{center}
    \textit{[Pantalla: ``Enter your email to continue'']}
    \end{center}

    \item Ingresa tu correo @udec.cl y haz clic en ``Send me a code''

    \item Recibir\'as un c\'odigo de verificaci\'on en tu email:

    \begin{center}
    \textit{[Pantalla: ``Enter the 6-digit code sent to your email'']}
    \end{center}

    \item Ingresa el c\'odigo de 6 d\'igitos

    \item Acceder\'as al dashboard de Grafana autom\'aticamente
\end{enumerate}

\begin{tipbox}
Solo necesitas verificar tu email la primera vez. Despu\'es, el acceso se recordar\'a por 24 horas.
\end{tipbox}

\subsection{Accesos Posteriores}

Una vez verificado tu email, los accesos posteriores son directos: simplemente abre \url{https://lautaro.elprobedor.com} y estar\'as en el dashboard.

% ============================================================
\section{Dashboard: Tiempo Real}
% ============================================================

Al acceder al sistema, ver\'as el dashboard principal ``Solar Monitor - Tiempo Real'', que muestra los datos actuales del inversor actualizados cada 5 segundos.

\subsection{Indicadores Principales}

En la parte superior encontrar\'as 4 indicadores tipo gauge:

\begin{table}[H]
\centering
\begin{tabular}{llp{8cm}}
\toprule
\textbf{Indicador} & \textbf{Unidad} & \textbf{Descripci\'on} \\
\midrule
Potencia AC & W & Potencia instant\'anea entregada a la red el\'ectrica \\
Temperatura & °C & Temperatura interior del inversor \\
Energ\'ia Diaria & kWh & Energ\'ia generada en el d\'ia actual \\
Energ\'ia Total & kWh & Energ\'ia total generada desde la instalaci\'on \\
\bottomrule
\end{tabular}
\end{table}

\subsection{Gr\'aficos de Tiempo}

Debajo de los indicadores encontrar\'as gr\'aficos de l\'nea con los datos de las \'ultimas 24 horas:

\begin{itemize}
    \item \textbf{Potencia AC}: Muestra la evoluci\'on de la potencia a lo largo del d\'ia. \'Util para identificar la curva de producci\'on solar.
    \item \textbf{Voltaje y Corriente DC}: Muestra el voltaje ($V_{pv}$) y la corriente ($I_{pv}$) del arreglo fotovoltaico.
    \item \textbf{Voltaje y Frecuencia AC}: Muestra el voltaje de red ($V_{ac}$) y la frecuencia ($F_{ac}$).
\end{itemize}

\subsection{Estado del Inversor}

En la secci\'on de estado encontrar\'as:

\begin{table}[H]
\centering
\begin{tabular}{lll}
\toprule
\textbf{Estado} & \textbf{Color} & \textbf{Significado} \\
\midrule
OK & Verde & Inversor operando normalmente \\
Standby & Amarillo & Inversor en espera (sin producci\'on) \\
Off & Gris & Inversor apagado \\
Fault & Rojo & Inversor en condici\'on de falla \\
\bottomrule
\end{tabular}
\caption{Estados del inversor}
\end{table}

\begin{warnbox}
Si el estado muestra ``Fault'' (rojo), revisar la secci\'on de diagn\'ostico o contactar al administrador.
\end{warnbox}

% ============================================================
\section{Dashboard: Hist\'orico}
% ============================================================

El dashboard ``Solar Monitor - Hist\'orico'' permite analizar datos pasados del inversor.

\subsection{Energ\'ia Diaria}

Muestra un gr\'afico de barras con la energ\'ia generada cada d\'ia del \'ultimo mes. Permite identificar d\'ias con baja producci\'on (nublados) y d\'ias con alta producci\'on (despejados).

\subsection{Energ\'ia Mensual}

Muestra la energ\'ia total generada cada mes del \'ultimo a\~no. \'Util para analizar la estacionalidad de la producci\'on solar.

\subsection{Temperatura vs Potencia}

Gr\'afico de dispersi\'on que relaciona la temperatura del inversor con la potencia generada. Permite observar la degradaci\'on t\'ermica del inversor a altas temperaturas.

\subsection{Eficiencia del Inversor}

Calcula la eficiencia del inversor como:

$$\eta = \frac{P_{ac}}{V_{pv} \times I_{pv}} \times 100\%$$

Donde:
\begin{itemize}
    \item $P_{ac}$: Potencia AC entregada a la red [W]
    \item $V_{pv}$: Voltaje DC del arreglo fotovoltaico [V]
    \item $I_{pv}$: Corriente DC del arreglo fotovoltaico [A]
\end{itemize}

% ============================================================
\section{Dashboard: Diagn\'ostico}
% ============================================================

El dashboard ``Solar Monitor - Diagn\'ostico'' muestra la salud del sistema de monitoreo.

\subsection{Indicadores Clave}

\begin{itemize}
    \item \textbf{\'Ultima Lectura Exitosa}: Verde si es reciente ($<30$ seg), amarillo si es antigua ($<5$ min), rojo si es muy antigua ($>5$ min)
    \item \textbf{Intervalo entre Lecturas}: Deber\'ia ser estable en $\sim$5 segundos
    \item \textbf{Errores de Comunicaci\'on}: N\'umero de errores Modbus por hora
    \item \textbf{Temperatura del Inversor}: Con l\'imite de alerta en 65 °C
    \item \textbf{Voltaje de Red}: Con bandas aceptables 220 V $\pm$ 10\%
    \item \textbf{Frecuencia de Red}: Con bandas aceptables 50 Hz $\pm$ 0.5 Hz
\end{itemize}

% ============================================================
\section{Dashboard: Acad\'emico}
% ============================================================

El dashboard ``Solar Monitor - Acad\'emico'' proporciona indicadores relevantes para an\'alisis acad\'emico.

\subsection{Horas de Sol Equivalentes}

Las horas de sol equivalentes (HSE) se calculan como:

$$HSE = \frac{E_{diaria}}{P_{nominal}} = \frac{\sum P_{ac} \times \Delta t}{4000}$$

Donde $P_{nominal} = 4000$ W es la potencia nominal del inversor.

\subsection{Energ\'ia Espec\'ifica}

La energ\'ia espec\'ifica se expresa en kWh/kWp:

$$Y_f = \frac{E_{diaria}}{P_{nominal}}$$

Donde $P_{nominal}$ corresponde a la potencia pico del arreglo FV en kWp.

\subsection{Performance Ratio (PR)}

El Performance Ratio es el indicador m\'as importante para evaluar la calidad de un sistema FV:

$$PR = \frac{Y_f}{Y_r} = \frac{E_{real}}{E_{teorica}} \times 100\%$$

Donde $Y_r$ es la irradiaci\'on en el plano del arreglo dividida por la irradiancia de referencia (1000 W/m$^2$).

\begin{infobox}
El c\'alculo del PR requiere datos de irradiancia solar. Si no se dispone de piran\'omero, se pueden utilizar datos de la DMC o APIs como Solcast/PVGIS para la ubicaci\'on de Concepci\'on ($-36.83$, $-73.04$).
\end{infobox}

\subsection{CO2 Evitado}

El inversor reporta el CO2 ahorrado acumulado. Este valor se calcula internamente por el inversor asumiendo un factor de emisi\'on de la red el\'ectrica chilena.

% ============================================================
\section{Exportar Datos}
% ============================================================

\subsection{Exportar CSV desde Grafana}

Para descargar datos en formato CSV:

\begin{enumerate}
    \item Hacer clic en el t\'itulo del panel que deseas exportar
    \item Seleccionar ``Inspect'' del men\'u
    \item Ir a la pesta\~na ``Data''
    \item Hacer clic en ``Download CSV''
\end{enumerate}

\begin{tipbox}
Para exportar datos de un rango de fechas espec\'ifico, usa el selector de tiempo en la esquina superior derecha de Grafana antes de exportar.
\end{tipbox}

\subsection{Exportar JSON desde Grafana}

Tambi\'en se puede exportar en formato JSON:

\begin{enumerate}
    \item Hacer clic en el t\'itulo del panel que deseas exportar
    \item Seleccionar ``Inspect'' del men\'u
    \item Ir a la pesta\~na ``Data''
    \item Hacer clic en ``Download JSON''
\end{enumerate}

\subsection{Consulta directa a TimescaleDB (avanzado)}

Para usuarios con acceso SSH al servidor lautaro, se pueden hacer consultas SQL directas a la base de datos:

\begin{lstlisting}[language=bash, caption={Ejemplos de consulta SQL directa}]
# Conectar a la base de datos
docker exec -it solar-monitor-timescaledb-1 psql -U solar solar_monitor

# Ultimos 10 registros en tiempo real
SELECT * FROM realtime ORDER BY time DESC LIMIT 10;

# Energia diaria de los ultimos 7 dias
SELECT bucket, avg_power_w, peak_power_w
FROM daily_energy
WHERE bucket > NOW() - INTERVAL '7 days'
ORDER BY bucket;

# Eventos de las ultimas 24 horas
SELECT time, event_type, event_value, severity
FROM events
WHERE time > NOW() - INTERVAL '24 hours'
ORDER BY time DESC;
\end{lstlisting}

\subsection{Usar datos en Python (avanzado)}

Para usuarios avanzados con acceso SSH al servidor, se puede consultar TimescaleDB directamente desde Python:

\begin{lstlisting}[language=Python, caption={Ejemplo: Consultar datos con Python}]
import pandas as pd
import psycopg2

# Conectar a TimescaleDB (requiere acceso SSH o red local)
conn = psycopg2.connect(
    host="192.168.1.145",
    port=5432,
    database="solar_monitor",
    user="grafana_reader",
    password="tu_password"
)

# Obtener datos de las ultimas 24 horas
query = """
SELECT time_bucket('1 minute', time) AS bucket,
       AVG(pac) AS avg_power_w
FROM fast_samples
WHERE time > NOW() - INTERVAL '24 hours'
GROUP BY bucket
ORDER BY bucket;
"""
df = pd.read_sql(query, conn)
conn.close()

# Calcular energia diaria
energy_kwh = df["avg_power_w"].sum() * 1/60 / 1000
print(f"Energia diaria: {energy_kwh:.2f} kWh")
\end{lstlisting}

\begin{warnbox}
El acceso directo a TimescaleDB requiere estar en la red local (192.168.1.145) o conectado v\'ia Tailscale (100.66.126.91). No est\'a accesible desde internet por seguridad.
\end{warnbox}

% ============================================================
\section{Variables Disponibles}
% ============================================================

Las variables que el sistema puede monitorear son:

\begin{table}[H]
\centering
\begin{tabular}{llll}
\toprule
\textbf{Variable} & \textbf{S\'imbolo} & \textbf{Unidad} & \textbf{Descripci\'on} \\
\midrule
Voltaje DC & $V_{pv}$ & V & Voltaje del arreglo FV \\
Corriente DC & $I_{pv}$ & A & Corriente del arreglo FV \\
Voltaje AC & $V_{ac}$ & V & Voltaje de red \\
Corriente AC & $I_{ac}$ & A & Corriente inyectada a la red \\
Frecuencia AC & $F_{ac}$ & Hz & Frecuencia de red \\
Potencia AC & $P_{ac}$ & W & Potencia entregada \\
Impedancia AC & $Z_{ac}$ & $\Omega$ & Impedancia de red \\
Temperatura & $T$ & °C & Temperatura del inversor \\
Energ\'ia total & $E_{total}$ & kWh & Energ\'ia acumulada \\
Energ\'ia diaria & $E_{daily}$ & kWh & Energ\'ia del d\'ia \\
Horas operaci\'on & $H_{total}$ & h & Horas totales de operaci\'on \\
CO2 ahorrado & $CO_2$ & kg & CO2 evitado acumulado \\
Estado & Status & --- & Estado del inversor \\
Estado red & Grid & --- & Estado de conexi\'on a red \\
\bottomrule
\end{tabular}
\caption{Variables disponibles del inversor}
\end{table}

% ============================================================
\section{Frecuentes Problemas}
% ============================================================

\subsection{No puedo acceder al sistema}

\begin{enumerate}
    \item Verificar que est\'as usando la URL correcta: \url{https://lautaro.elprobedor.com}
    \item Verificar que tu email sea @udec.cl (o @ing.udec.cl, @dci.udec.cl)
    \item Revisar la carpeta de spam para el c\'odigo de verificaci\'on
    \item Si persiste, contactar al administrador: daniel.ruminot.moscoso@gmail.com
\end{enumerate}

\subsection{Los datos no se actualizan}

\begin{enumerate}
    \item Verificar la conexi\'on a internet
    \item Hacer clic en el bot\'on ``Refresh'' en Grafana (esquina superior derecha)
    \item Verificar que el inversor est\'e encendido (no es de noche ni est\'a en standby)
    \item Revisar el dashboard de Diagn\'ostico para ver el estado del sistema
\end{enumerate}

\subsection{Los gr\'aficos est\'an vac\'ios}

\begin{enumerate}
    \item Verificar el rango de tiempo seleccionado (esquina superior derecha)
    \item El inversor no produce energ\'ia de noche --- probar con ``Last 7 days'' en lugar de ``Last 30 minutes''
    \item Si persiste de d\'ia, contactar al administrador
\end{enumerate}

\subsection{Error ``Too many requests''}

El sistema tiene un l\'imite de solicitudes. Esperar 1 minuto y reintentar. Si el problema persiste, contactar al administrador.

% ============================================================
\section{Glosario}
% ============================================================

\begin{table}[H]
\centering
\begin{tabular}{ll}
\toprule
\textbf{T\'ermino} & \textbf{Definici\'on} \\
\midrule
FV & Fotovoltaico --- conversi\'on de luz solar en electricidad \\
Inversor & Dispositivo que convierte DC en AC \\
AC & Corriente alterna (red el\'ectrica) \\
DC & Corriente continua (paneles solares) \\
Modbus & Protocolo de comunicaci\'on industrial \\
RS485 & Est\'andar de comunicaci\'on serial diferencial \\
kWh & Kilovatio-hora --- unidad de energ\'ia \\
kWp & Kilovatio pico --- potencia nominal del arreglo FV \\
PR & Performance Ratio --- eficiencia real vs.\ te\'orica \\
HSE & Horas de sol equivalentes \\
Grafana & Plataforma de visualizaci\'on de datos \\
TimescaleDB & Base de datos de series temporales \\
\bottomrule
\end{tabular}
\end{table}

% ============================================================
\section{Contacto}
% ============================================================

Para consultas o problemas con el sistema:

\begin{itemize}
    \item \textbf{Administrador}: Daniel Ruminot Moscoso
    \item \textbf{Email}: daniel.ruminot.moscoso@gmail.com
    \item \textbf{Departamento}: Ingenier\'ia El\'ectrica, Universidad de Concepci\'on
\end{itemize}

\end{document}